%% ============================================================
%% CHAPTER 6 — Context Compaction for Device-Hosted Loops
%%                                        (budget: ~7,000 words)
%% Role: answers RQ2. The D2 study grown to chapter length.
%% ============================================================
\chapter{Context Compaction for Device-Hosted Loops}
\label{ch:compaction}

\section{The Within-Turn Re-Upload Problem}
\label{sec:comp-problem}
% The code-level finding: the loop deep-copies and re-uploads the full
% message array on every tool iteration; within-turn uplink is O(k^2) in
% iteration count. Derivation + a measured motivating trace (bytes per
% iteration for a 10-iteration task). Why prompt caching does not help
% the device: bytes still leave the radio; RAM still holds the array.

\section{The Compaction Ladder}
\label{sec:comp-ladder}
% Design of the four levels, each a firmware switch:
% L0 full append (baseline); L1 per-tool output budgets with head/tail
% truncation; L2 deterministic iteration digests (tool_use/tool_result
% pairing preserved — API constraint — results shrunk to one status line,
% interim text dropped); L3 placeholder masking (the on-device
% reproduction of observation masking \citep{lindenbauer2025complexity}).
% Information-preservation argument for L2: what a next-step decision
% actually needs.

\section{Implementation}
\label{sec:comp-impl}
% context_compactor module (~80 LoC), per-tool budget field in the tool
% registry (~30 LoC), the is_error repair to tool results. Where the
% compactor runs in the loop; why in-place cJSON rewriting is safe.

\section{Evaluation}
\label{sec:comp-eval}
% Per level: bytes-per-turn curve (the O(k^2)->O(k) figure), tokens,
% per-iteration latency, PSRAM high-water, success rate on the task grid
% (stratified by horizon and recall requirement). The key trade-off
% figure: success vs uplink across L0–L3. Model-tier interaction: which
% level weaker models tolerate.

\section{Relation to State-Folding Approaches}
\label{sec:comp-related}
% Position L1/L2 as the deterministic rungs below MEM1/state-folding
% (L3+ requires the model in the loop); what the ladder buys without any
% model cooperation; when it stops being enough.

\section{Summary}
\label{sec:comp-summary}
% One paragraph: the recommendation (default level, when to deviate).
